% ============================================================
\subsubsection{Ordering and Row-Relative Logic}
% ============================================================

Some results depend on the order of the rows or on values in other rows. Ordering operations sort or rank the frame, and row-relative operations read a neighbor or a running window. These are the pandas counterpart of SQL window functions, and as with a SQL window the order must be stated, since a frame carries no inherent one.

\paragraph{Sorting.} \texttt{sort\_values(by, ascending, na\_position)} orders rows by one or more columns, with \texttt{by} taking a list for a multi-key sort and \texttt{ascending} a matching list of directions. \texttt{na\_position} places \texttt{NaN} at the start or the end. \texttt{sort\_index} orders by the index labels instead.

\begin{lstlisting}[style=python, caption={Multi-key sort with a direction per key}]
df.sort_values(['dept', 'salary'], ascending=[True, False])
\end{lstlisting}

\vspace{-1ex}

\paragraph{Extremes and top-N.} \texttt{nlargest(n, col)} and \texttt{nsmallest(n, col)} return the top or bottom rows by a column and cost less than a full sort for a handful of rows. \texttt{idxmax} and \texttt{idxmin} return the index label of the extreme value, an argmax, which \texttt{loc} then reads as a row. Both break ties by first occurrence.

\begin{lstlisting}[style=python, caption={Top rows and the row of the maximum}]
df.nlargest(3, 'salary')
df.loc[df['salary'].idxmax()]      # the single row holding the max
\end{lstlisting}

\vspace{-1ex}

\paragraph{The Nth distinct value.} The second-highest and Nth-highest questions read the value sitting at a chosen rank. Deduplicating the values first keeps ties from consuming a rank, and the result resolves to a missing value when the data holds fewer than \texttt{N} distinct entries, which is the required answer to an absent Nth place.

\begin{lstlisting}[style=python, caption={Nth-highest distinct value, NaN when it does not exist}]
def nth_highest(s, n):
    u = s.drop_duplicates().nlargest(n)
    return u.iloc[-1] if len(u) >= n else np.nan
\end{lstlisting}

\vspace{-1ex}

\paragraph{Ranking.} \texttt{rank(method, ascending)} numbers rows by value, with the tie policy set by \texttt{method}: \texttt{dense}, \texttt{min}, and \texttt{first} correspond to the SQL \texttt{DENSE\_RANK}, \texttt{RANK}, and \texttt{ROW\_NUMBER}. Applied after a \texttt{groupby} it ranks within each group, the top-$N$-per-group tool. Sorting and then \texttt{groupby(...).head(n)} reaches the same result by keeping the first \texttt{n} rows of each group. Passing \texttt{pct=True} scales the ranks to the interval $(0, 1]$, giving a percent rank.

\begin{lstlisting}[style=python, caption={Rank within a group, and top-N per group by slicing}]
df['rk'] = df.groupby('dept')['salary'].rank(method='dense', ascending=False)

(df.sort_values('salary', ascending=False)
   .groupby('dept').head(1))        # highest-paid row per dept
\end{lstlisting}

\vspace{-1ex}

\paragraph{Position within a group.} \texttt{groupby(...).cumcount()} numbers the rows of each group by their current order from 0 and resets at every group edge, the within-group \texttt{ROW\_NUMBER} carrying no value to rank on. \texttt{ngroup()} labels the groups themselves \texttt{0, 1, 2} in order of appearance. Neither reads a column, so both rest on a prior sort the way \texttt{rank} does.

\begin{lstlisting}[style=python, caption={A 0-based sequence number within each group}]
df['seq'] = df.sort_values('ts').groupby('user').cumcount()
\end{lstlisting}

\vspace{-1ex}

\begin{keybox}[Row-relative operations need an explicit order]
\texttt{shift}, \texttt{diff}, \texttt{rank}, \texttt{rolling}, and the cumulative methods all read the current row order and nothing else. Without a preceding \texttt{sort\_values} the previous row is whatever order the rows arrived in. This is the \texttt{ORDER BY} inside a SQL window, and a per-group version needs the operation applied through \texttt{groupby} so it resets at each group boundary rather than bleeding across.
\end{keybox}

\vspace{-1ex}

\paragraph{Neighbor access with shift and diff.} \texttt{shift(n)} moves a column down \texttt{n} positions and fills the vacated top with \texttt{NaN}, so \texttt{shift(1)} sets the previous row's value beside the current one, the \texttt{LAG}. \texttt{shift(-1)} looks ahead, the \texttt{LEAD}. \texttt{diff()} equals \texttt{current - shift(1)}, so \texttt{diff() > 0} tests whether a value rose. The first row of an ordering has no predecessor, so its shifted value is \texttt{NaN} and every comparison against it is \texttt{False}, which drops it with no explicit guard. Applying the shift through \texttt{groupby} resets it at each group edge. The one-period fractional change, \texttt{close / close.shift(1) - 1}, is packaged as \texttt{pct\_change}, the simple return.

\begin{lstlisting}[style=python, caption={Per-group return against the previous row}]
df = df.sort_values(['ticker', 'date'])

prev = df.groupby('ticker')['close'].shift(1)
df['ret'] = df['close'] / prev - 1               # explicit, via shift
df['ret'] = df.groupby('ticker')['close'].pct_change()   # same return
\end{lstlisting}

\vspace{-2.5ex}

\paragraph{Running and rolling windows.} A cumulative method such as \texttt{cumsum} or \texttt{cummax} runs from the start up to the current row, giving a running total. \texttt{rolling(k)} fixes a window of \texttt{k} consecutive rows for a moving statistic, and \texttt{expanding} grows the window from the start. \texttt{min\_periods} sets how many observations a window needs before it returns a value in place of \texttt{NaN}. Wrapping these in \texttt{groupby(...).transform(...)} keeps the result aligned to the original rows while resetting each group.

\begin{lstlisting}[style=python, caption={A per-group moving average aligned to the rows}]
df['ma20'] = (df.groupby('ticker')['close']
                .transform(lambda s: s.rolling(20).mean()))
\end{lstlisting}

\vspace{-1ex}

\texttt{rolling} accepts any reduction, so \texttt{.rolling(20).std()} is a moving volatility, and a moving correlation between two aligned series would be \texttt{.rolling(60).corr(other)}. \texttt{ewm} weights recent observations more heavily under an exponential decay set by \texttt{span}, \texttt{halflife}, or \texttt{alpha}, giving the exponentially weighted mean and variance used in volatility estimates. Compounding runs through \texttt{cumprod}, so \texttt{(1 + ret).cumprod() - 1} accumulates a return series.

\begin{lstlisting}[style=python, caption={Moving volatility, exponential mean, and compounded return}]
g = df.groupby('ticker')
df['vol20']   = g['ret'].transform(lambda s: s.rolling(20).std())
df['ewma20']  = g['close'].transform(lambda s: s.ewm(span=20).mean())
df['cum_ret'] = g['ret'].transform(lambda s: (1 + s).cumprod() - 1)
\end{lstlisting}

\vspace{-1ex}

\paragraph{Time-based and custom windows.} An integer \texttt{rolling(k)} counts \texttt{k} rows and ignores their timestamps. An offset string, \texttt{rolling('7D')}, counts a span of calendar time instead and needs a \texttt{DatetimeIndex} or an \texttt{on=} column naming the time field. The two coincide only when the rows sit one period apart with no gap between them. A window takes any reduction, so \texttt{sum}, \texttt{min}, \texttt{max}, \texttt{median}, and \texttt{quantile} all apply, alongside \texttt{center=True} for a window centered on each row. \texttt{rolling(k).apply(f, raw=True)} runs an arbitrary \texttt{f} over each window as a numpy array, the route for a statistic with no built-in form.

\begin{lstlisting}[style=python, caption={A calendar window, a custom window, and a centered median}]
df = df.set_index('date').sort_index()

df['roll7'] = df['amount'].rolling('7D').sum()     # seven calendar days back
df['rng10'] = df['amount'].rolling(10).apply(lambda w: w.max() - w.min(), raw=True)
df['med5']  = df['amount'].rolling(5, center=True).median()
\end{lstlisting}

\vspace{2ex}

\begin{keybox}[A row window is not a day window]
\texttt{rolling(7)} spans seven rows, which equals seven days only when every day is present. On a series with missing dates it reaches further back in time than seven days. Two fixes hold the window to calendar length: reindex to a complete date range so each day is a row, or use the offset form \texttt{rolling('7D')}, which measures the span directly and covers the rows whose timestamps fall within seven days back of the current one.
\end{keybox}

\vspace{-1ex}

\paragraph{Time indexing and resampling.} A \texttt{DatetimeIndex} unlocks time-aware operations the ordinary index lacks. Partial-string slicing selects by calendar span, so \texttt{df.loc['2023']} takes the whole year and \texttt{df.loc['2023-01':'2023-06']} a range of months. \texttt{resample} is a groupby keyed by a time frequency: it bins the rows into calendar buckets and reduces each, downsampling a series to a coarser frequency or upsampling to a finer one that then needs filling. Frequency aliases name the bucket, among them \texttt{'B'} business day, \texttt{'W'} week, \texttt{'ME'} month-end, and \texttt{'QE'} quarter-end, the last two written \texttt{'M'} and \texttt{'Q'} before pandas 2.2.

\begin{lstlisting}[style=python, caption={Downsample, build OHLC bars, and align to business days}]
px = px.set_index('date').sort_index()

px['close'].resample('ME').last()      # month-end close (downsample)
px['close'].resample('W').ohlc()       # weekly open-high-low-close bars
px['close'].resample('B').ffill()      # to business days, carry last forward

px.groupby('ticker')['close'].resample('ME').last()   # per entity
\end{lstlisting}
